\chapter{Computação e Música}
\label{cap:computacao_musical}

Este capítulo apresenta a base teórica que serviu de apoio para o desenvolvimento da solução computacional do Aerochord. Na Seção~\ref{sec:fundamentos_computacao_musical}, discutem-se os fundamentos da computação musical, incluindo representação digital do som, síntese sonora e processamento de áudio digital. Em seguida, a Seção~\ref{sec:protocolo_midi} detalha o protocolo MIDI (Interface Digital para Instrumentos Musicais, do inglês \textit{Musical Instrument Digital Interface}), apresentando seu propósito, funcionamento e as mensagens mais relevantes para o controle expressivo do sistema. Por fim, na Seção~\ref{sec:teoria_musical}, revisam-se os principais conceitos de teoria musical aplicados ao mapeamento gestual — escala cromática, oitavas, articulações (notas curtas e destacadas, \textit{staccato}; ou notas fluidas e ligadas, \textit{legato}) e efeitos como modulação de afinação e frequência (\textit{pitch bend} e \textit{vibrato}), sustentação (\textit{sustain}) e dinâmica — que fundamentaram o \textit{design} do instrumento.

\section{Fundamentos da Computação Musical}
\label{sec:fundamentos_computacao_musical}

A computação musical concentra-se na aplicação de algoritmos e modelos computacionais para criar, analisar e modificar sons de maneira flexível e precisa. Conforme destacado por Miranda (\citeyear{Miranda2002}), esse campo abrange desde a geração programada de timbres até a modelagem de estruturas musicais complexas e sistemas de execução (\textit{performance}) interativa. Dodge e Jerse (\citeyear{Dodge1997}) enfatizam o papel do computador na realização de tarefas sonoras avançadas, enquanto Moore (\citeyear{Moore1990}) evidencia a importância da representação digital do som, permitindo seu armazenamento e processamento eficientes.

O processo inicia-se com a conversão de um sinal analógico em dados numéricos, por meio de dois estágios principais: a amostragem, que mede a amplitude do sinal em instantes regulares, e a quantização, que mapeia essas medições para valores discretos \cite{Rossing2002}. Essa sequência de valores numéricos serve de base para todas as operações subsequentes.

A síntese sonora gera sons por meio de modelos computacionais em vez de gravações acústicas. Na síntese aditiva, ondas senoidais puras são somadas para formar timbres; na síntese subtrativa, sinais ricos em harmônicos são filtrados para moldar o espectro sonoro. A modulação de frequência (FM) altera dinamicamente a frequência de uma portadora por uma moduladora, produzindo timbres expressivos, enquanto a síntese granular emprega pequenos fragmentos de áudio (os ``grãos'') para criar texturas sonoras singulares \cite{Manning1993}.

O processamento de áudio digital manipula gravações pré-existentes por meio de técnicas como filtragem, equalização, compressão, reverberação e aplicação de efeitos especiais, aprimorando a qualidade sonora e possibilitando transformações criativas \cite{Zolzer2011}. Essas operações são largamente empregadas em Estações de Trabalho de Áudio Digital (\textit{Digital Audio Workstations} - DAWs) para modelar o som conforme a intenção artística.

As aplicações da computação musical incluem a criação de instrumentos digitais equipados com sensores e controladores, a composição assistida por computador que apoia o desenvolvimento de ideias melódicas e estruturais, e sistemas interativos que respondem em tempo real aos gestos do intérprete (\textit{performer}). Adicionalmente, métodos de análise e recuperação de informação musical permitem indexar e buscar repertórios com base em atributos como melodia, ritmo e timbre \cite{Cope2001}.

No contexto deste trabalho, esses fundamentos habilitaram a captação e interpretação de gestos corporais e seu mapeamento para mensagens MIDI coerentes. A etapa de quantização assegurou precisão na conversão dos sinais de movimento, enquanto a síntese e o processamento digital compõem o ambiente sonoro no \textit{software} receptor. Ao integrar esses elementos, o hiperinstrumento desenvolvido amplia as possibilidades expressivas do intérprete, tornando seu corpo parte ativa da criação musical.

\section{O Protocolo MIDI (Musical Instrument Digital Interface)}
\label{sec:protocolo_midi}

Desde seu surgimento em meados de 1983, o protocolo MIDI estabeleceu-se como padrão onipresente para comunicação de eventos musicais entre controladores, sintetizadores e sequenciadores, permitindo que equipamentos de fabricantes distintos interagissem de forma consistente e interoperável \cite{SMC2023MIDI,MIDI1DetailedSpec}. Essa universalidade viabilizou a criação de um ecossistema diversificado de equipamentos físicos (\textit{hardware}) e lógicos (\textit{software}), moldando práticas de produção, \textit{performance} e composição musical ao longo de décadas. Porém, à medida que as demandas por expressividade mais refinada, controle de alta precisão e interação em tempo real evoluíram, especialmente em contextos de performance gestual, tornaram-se evidentes limitações do MIDI 1.0, tais como resolução restrita de parâmetros, comunicação predominantemente unidirecional e ausência de mecanismos de autoconfiguração \cite{MIDI2Overview2023}.

Em cenários de interfaces musicais gestuais de baixa latência e alta expressividade, como o projeto Aerochord, essas restrições tornam-se especialmente críticas: o mapeamento fino de gestos a parâmetros sonoros exige curvas de resposta suaves, confirmação de capacidades do dispositivo receptor e sincronização precisa de eventos. Foi essa necessidade que motivou a adoção do MIDI 2.0 no projeto, cuja arquitetura revisada introduz inovações concebidas para atender a tais exigências sem abandonar a compatibilidade retroativa com o ecossistema existente. Nesta seção, são apresentados inicialmente os fundamentos e as limitações do protocolo MIDI 1.0. Em seguida, discute-se a arquitetura técnica e as novas capacidades do MIDI 2.0, demonstrando como elas fundamentaram a escolha desse protocolo como pilar tecnológico do Aerochord.

\subsection{Fundamentos e Limitações do MIDI 1.0}
O protocolo MIDI 1.0 opera como um sistema de transmissão de instruções musicais, transportando mensagens que indicam eventos sem veicular o áudio propriamente dito \cite{SMC2023MIDI,MIDI1DetailedSpec}. Os eventos mais comuns incluem ativação e desativação de notas (\textit{Note On} e \textit{Note Off}), intensidade do toque (\textit{Velocity}), mudanças de parâmetros de som (\textit{Control Change}), troca de instrumentos (\textit{Program Change}), dobra de afinação (\textit{Pitch Bend}) e mensagens exclusivas de fabricantes (\textit{System Exclusive} ou \textit{SysEx}). Cada mensagem é composta por um byte de estado (com bit alto indicando tipo de mensagem e canal) seguido de um ou dois bytes de dados (com 7 bits significativos), refletindo a concepção original de comunicação serial a 31,25 kbps. Essa taxa de transferência definia limites para o número de eventos simultâneos que poderiam ser enviados sem comprometer significativamente a latência ou gerar congestionamento na linha.

Ao unificar equipamentos de fabricantes distintos, o MIDI 1.0 estabeleceu um ecossistema interoperável, permitindo que controladores e DAWs compartilhassem mensagens comuns. Essa compatibilidade retroativa sustenta uma vasta base instalada de dispositivos e aplicou-se também ao Aerochord, que implementou a capacidade de enviar mensagens em modo legado a sintetizadores sem suporte a MIDI 2.0, assegurando que usuários com equipamentos existentes possam integrar o sistema sem barreiras.

Entretanto, para aplicações que demandam expressividade refinada e baixa latência a concepção original do MIDI 1.0 revela limitações críticas. A resolução de 7 bits em \textit{Control Changes} e \textit{Velocity} implica apenas 128 passos discretos, frequentemente resultando em saltos audíveis em variações contínuas; o \textit{Pitch Bend} afeta todas as notas de um mesmo canal simultaneamente, impedindo modulações individuais em acordes polifônicos; a comunicação é primordialmente unidirecional; o número restrito de 16 canais por porta limita distinções granulares; e a ausência de marcas de tempo (\textit{timestamps}) embutidas impede compensações automáticas de variações no tempo de chegada dos dados (\textit{jitter})\footnote{Variação no atraso da transmissão de dados sobre uma rede ou barramento, causando irregularidades e possíveis problemas de sincronização.}. Esses fatores se tornam especialmente significativos no Aerochord, em que a fluidez do mapeamento de gestos exigiu curvas suaves e sincronização precisa que ultrapassam o escopo do padrão legado \cite{MIDI2Overview2023}.

\subsection{Arquitetura Técnica do MIDI 2.0}
O surgimento do MIDI 2.0 representa uma resposta direta às limitações identificadas no protocolo legado, oferecendo compatibilidade retroativa ao introduzir mecanismos bidirecionais que permitem autoconfiguração e descoberta de capacidades entre instrumentos \cite{MIDI2Overview2023}. Essa evolução foi motivada pela necessidade de suportar demandas de expressividade refinada em tempo real, onde a fluidez gestual é um requisito central.

No cerne da nova arquitetura estão os Pacotes MIDI Universais (\textit{Universal MIDI Packets} - UMP), que redefinem o formato de mensagens para 32, 64 ou 128 bits, com campos que incluem tipo de mensagem, grupo, canal/comando e, opcionalmente, \textit{timestamp} embutido \cite{UMPProtocol2023}. Essa estrutura ampliada possibilita transportar dados de alta resolução e atuar com controladores por nota individual (\textit{per-note controllers})\footnote{Controladores que permitem ajustar parâmetros como volume, afinação e timbre de cada nota individualmente, diferentemente dos controles comuns que afetam todo o canal.}. No Aerochord, cada evento gestual é marcado com um \textit{timestamp} no momento exato da captura, propagado internamente por todo o \textit{pipeline} até o envio do pacote MIDI. Esse \textit{timestamp} é usado para medir a latência \textit{end-to-end} do sistema (percentis P50/P95) — ele não é serializado como mensagem de \textit{Jitter-Reduction Timestamp} dentro do próprio pacote UMP transmitido ao dispositivo receptor, ficando restrito à instrumentação interna do Aerochord.

A descoberta de capacidades passa a ser realizada por meio da Consulta de Capacidades MIDI (\textit{MIDI Capability Inquiry} - MIDI-CI), em que um dispositivo envia mensagens de busca e recebe detalhes sobre o suporte a MIDI 2.0 e propriedades configuráveis \cite{MIDICI2023}. Esse protocolo de negociação inicial (\textit{handshake}) transforma a comunicação "cega" em uma negociação inteligente: o Aerochord identifica automaticamente se o receptor aceita mensagens UMP, dispensando configurações manuais extensas.

Em sequência, os Perfis (\textit{Profiles}) e a Troca de Propriedades (\textit{Property Exchange}) permitem compartilhar conjuntos padronizados de regras. Os perfis definem como interpretar mensagens em contextos específicos (ex: mapeamentos para \textit{per-note controllers}), enquanto o \textit{Property Exchange} possibilita consultar e ajustar dinamicamente parâmetros remotos, como curvas de resposta, sem sair da sessão \cite{MIDICIProfiles2023,MIDIPropertyExchange2023}. 

A arquitetura também prevê flexibilidade de transporte: pacotes UMP podem trafegar via USB, \textit{Bluetooth} ou rede, ao mesmo tempo em que existe uma contingência (\textit{fallback}) para encapsular mensagens em conexões MIDI 1.0 legadas \cite{UMPProtocol2023}. Essa adaptabilidade garante operação em múltiplas plataformas, mantendo a funcionalidade básica mesmo em dispositivos antigos.

Por fim, a especificação prevê até 16 grupos independentes, cada um com 16 canais, permitindo separar fluxos de controle distintos. Como o Aerochord, em sua versão atual, é um instrumento monofônico de voz única — em que notas, volume, sustentação, timbre, vibrato e \textit{pitch bend} moldam o mesmo som —, todos os eventos das duas mãos são roteados para um único grupo e canal. Essa escolha garante que os macrocontroles globais (volume e sustentação, operados pela mão de controle) e as modulações por nota (operadas pela mão de articulação) incidam efetivamente sobre as notas em execução, uma vez que, na especificação, cada grupo constitui uma porta lógica de 16 canais independente. Os grupos adicionais permanecem reservados para a evolução polifônica e multivoz do sistema, discutida entre os trabalhos futuros (Capítulo~\ref{chap:consideracoes_finais}).

\subsection{Novas Capacidades Expressivas do MIDI 2.0}
Ao redefinir o formato e introduzir mecanismos de maior resolução, o MIDI 2.0 abre espaço para ganhos expressivos inimagináveis no padrão legado, os quais se refletiram diretamente no Aerochord \cite{MIDI2Overview2023,UMPProtocol2023}.

A alta resolução representa um avanço fundamental: valores de \textit{Velocity} passam a ter 16 bits e \textit{Control Changes} podem alcançar 32 bits, eliminando saltos audíveis. No Aerochord, onde pequenas alterações de gesto modulam contínua e simultaneamente vários efeitos, essa granularidade estendida assegurou que o mapeamento refletisse nuances sutis sem artefatos \cite{UMPProtocol2023}.

Outra capacidade expressiva é o suporte a controle por nota individual. Diferentemente do MIDI 1.0, o MIDI 2.0 permite enviar mensagens de afinação, modulação ou pressão após o toque (\textit{aftertouch}) para notas específicas dentro de um acorde. Esse recurso viabilizou que o Aerochord aplicasse afinação apenas a uma nota sustentada (\textit{UMP MT4 Opcode 0x6}), enquanto as outras permanecem estáveis, enriquecendo texturas polifônicas de modo antes inviável.

O transporte de blocos estruturados de dados torna-se mais robusto, permitindo enviar configurações e predefinições completas (\textit{presets}) de sintetizadores em formato padronizado. A padronização desses mecanismos facilita a portabilidade e a experimentação colaborativa, integrando-se à calibração dinâmica em tempo real \cite{MIDIPropertyExchange2023}. 

\subsection{Persistência de Performances e Arquivos}
Registrar performances gestuais com fidelidade é fundamental para edição e análise detalhada. Em instrumentos como o Aerochord, cada gesto traduzido em MIDI 2.0 contém dados de alta resolução e \textit{timestamps} que documentam a evolução temporal exata dos controles.

O formato Arquivo de Clipe MIDI (\textit{MIDI Clip File}), baseado em UMP, oferece suporte à gravação de sequências contendo mensagens de alta resolução e \textit{timestamps} embutidos \cite{MIDIClipFile2023}. Em complemento, o Arquivo Contêiner (\textit{Container File}) permite agrupar múltiplos clipes e metadados, formando um pacote coeso. 

No contexto de avaliação do Aerochord, a persistência de eventos (implementada através do modo de avaliação analítica, \textit{eval mode}, exportando para arquivos tabulares) e os registros (\textit{logs}) precisos de eventos de ponta a ponta (\textit{end-to-end}) desempenharam um papel vital na medição de latência e de precisão do sistema.

\subsection{Justificativa Final para a Arquitetura do Aerochord}
Em síntese, as capacidades do MIDI 2.0 — alta resolução, controle por nota individual, autoconfiguração via MIDI-CI e \textit{timestamps} — atenderam diretamente aos requisitos de expressividade e robustez que nortearam o Aerochord \cite{UMPProtocol2023,MIDICI2023,MIDIPropertyExchange2023,MIDIClipFile2023}.

Em termos de aplicação concreta, a distância entre os dedos pôde controlar filtros sonoros suavemente sem saltos audíveis. Ao sustentar acordes, gestos de vibração geraram modulação independente em notas específicas. Além disso, ao conectar-se internamente, o Aerochord integrou uma porta virtual em modo MIDI 2.0, configurando parâmetros (\textit{setup}) de forma transparente.

O suporte a MIDI 2.0 não foi meramente um recurso opcional, mas o elemento central para viabilizar o controle gestual de alta precisão alcançado. A adoção do padrão potencializou a missão de democratização, oferecendo uma exploração fina com resposta sonora refinada.

\section{Teoria Musical Aplicada}
\label{sec:teoria_musical}

Esta seção apresenta os fundamentos da teoria musical que serviram de base para o mapeamento gestual no sistema Aerochord. Serão abordados os conceitos de escala cromática e oitavas \cite{benward2014music}. Em seguida, noções de articulação, como \textit{staccato} e \textit{legato} \cite{duckworth2015creative}. Por fim, revisam-se efeitos como \textit{vibrato, pitch bend, sustain} e dinâmica, implementados a partir das variações nos gestos captados \cite{Dodge1997}.

\subsection{Escala Cromática e Oitavas}
\label{sec:teoria_musical_escalas}

A escala cromática consiste na divisão da oitava em doze semitons iguais, representando a menor unidade de distância entre duas notas no sistema temperado. Cada semitom equivale a um aumento fixo de frequência, criando um sistema padronizado \cite{benward2014music}.

No contexto do MIDI, cada nota é representada por um número inteiro de 0 a 127. A nota Dó central (C4 ou MIDI 60) serve como referência, onde aumentos de valor representam semitons sucessivos. As oitavas são formadas por doze semitons (relação de frequência 2:1), percebidas como uma repetição da mesma nota em um registro mais agudo \cite{duckworth2015creative}. Essa percepção invariável foi explorada no Aerochord por meio dos gestos de seleção.

Do ponto de vista do \textit{design} gestual, o mapeamento suporta usuários destros e canhotos sem qualquer configuração prévia: os papéis de cada mão (controle vs.\ articulação) são atribuídos pela ordem de entrada em quadro, e não pela lateralidade anatômica do usuário (ver Seção~\ref{sec:captura_deteccao}). A altura vertical do quadro de visão da câmera foi dividida em doze zonas discretas (semitons) ativadas pela \textbf{mão de articulação} (a segunda mão a entrar em quadro). Simultaneamente, movimentos verticais macro da \textbf{mão de controle} (a primeira a entrar) determinavam a transição entre oitavas. Isso possibilitou ao usuário navegar por registros diferentes de forma ergonômica, escolhendo livremente qual mão usar para cada papel.

\subsection{Articulações}
\label{sec:teoria_musical_articulacoes}

As articulações musicais determinam como as notas se conectam. O \textit{staccato} indica notas breves e destacadas. No Aerochord, esse efeito é reproduzido por movimentos rápidos e secos, como um fechamento súbito dos dedos seguido de uma liberação imediata, gerando notas curtas \cite{duckworth2015creative}.

Por outro lado, o \textit{legato} indica uma execução fluida. Em termos técnicos, significa a emissão de um \textit{Note On} antes do \textit{Note Off} da nota anterior, criando sobreposição. No Aerochord, isso foi interpretado permitindo o deslize da mão de articulação fechada entre as zonas verticais de nota. Essa abordagem exigiu lógica computacional precisa para gerir as fronteiras, utilizando zonas mortas (\textit{dead zones} - áreas de tolerância de entrada) e configurações de estabilização do sinal para evitar ativações duplas indesejadas \cite{benward2014music}.

\subsection{Efeitos de expressividade}
\label{sec:teoria_musical_efeitos}

A expressividade musical reflete a resposta a nuances da execução. O \textit{vibrato} é a modulação periódica na frequência da nota. No MIDI 2.0, seu controle opera com alta granularidade \cite{midi_org}. No Aerochord, esse efeito foi gerado a partir de pequenos movimentos oscilatórios laterais da mão de articulação durante a ativação da nota.

O \textit{pitch bend} altera a afinação de forma contínua. Tratado como \textit{Per-Note Pitch Bend} (atributo por nota) em MIDI 2.0, o sistema interpretou deslocamentos angulares de inclinação da mão de articulação para realizar transições melódicas fluídas, mapeando a rotação diretamente para o valor de afinação da nota em andamento \cite{midi_org}.

O \textit{sustain} refere-se à capacidade de manter o som após o término da ação física. A manutenção da \textbf{mão de controle} aberta (em repouso vertical temporário) foi mapeada para acionar e desativar o \textit{sustain}, garantindo que as notas ativadas pela \textbf{mão de articulação} pudessem permanecer soando de maneira análoga ao pedal de sustentação de um piano.

A dinâmica diz respeito à intensidade sonora (comumente tratada na música tradicional sob termos como \textit{piano}, \textit{forte} e \textit{crescendo}). No protocolo MIDI, a dinâmica do ataque da nota é controlada pelo \textit{velocity}. No projeto, a velocidade de fechamento dos dedos da mão de articulação (distância mensurada em janelas de tempo) foi calculada pelo \textit{software} e mapeada para a intensidade da nota, permitindo ataques percussivos suaves ou agressivos conforme a intenção do usuário \cite{benward2014music}.

\section{Considerações sobre o Capítulo}
\label{sec:consideracoes_cap3}

Ao longo deste capítulo, foram apresentados os fundamentos da computação musical e explorado o protocolo MIDI. Demonstrou-se como o MIDI 2.0 reverteu obstáculos de baixa resolução por meio de pacotes universais (UMP) e consultas automáticas (MIDI-CI).

A compreensão dessas especificações técnicas constituiu o alicerce sobre o qual a arquitetura de \textit{software} do Aerochord se apoiou. A capacidade de endereçar dados de alta resolução e controlar a afinação (\textit{pitch bend}) de cada nota individualmente foi o elemento-chave que viabilizou o mapeamento fluido implementado.

Com os conceitos de IHC, teoria musical e protocolo MIDI firmemente estabelecidos, o Capítulo~\ref{cap:analise_estado_arte_interfaces_musicais_gestuais} aprofundará a revisão da literatura, examinando soluções existentes para delimitar exatamente o cenário onde a proposta de solução (detalhada no Capítulo~\ref{chap:desenvolvimento_proposta}) foi implementada.